\documentclass[11pt]{article}
\newcommand{\documenttype}{Curriculum vitae}
\usepackage[a4paper,margin=0.85in]{geometry}
\usepackage[T1]{fontenc}
\usepackage[utf8]{inputenc}
\usepackage{lmodern,amsmath,amssymb,microtype,xcolor,enumitem,tabularx,titlesec,fancyhdr,hyperref}
\definecolor{accent}{HTML}{315B8A}
\definecolor{link}{HTML}{5B4F7A}
\hypersetup{colorlinks=true,urlcolor=link,linkcolor=link,citecolor=link,pdfauthor={Phurailatpam Hemantakumar}}
\setlength{\parindent}{0pt}\setlength{\parskip}{5pt}\setlength{\headheight}{14pt}\setlength{\emergencystretch}{2em}
\titleformat{\section}{\large\bfseries\color{accent}}{}{0pt}{}
\titlespacing*{\section}{0pt}{9pt}{4pt}
\setlist[itemize]{leftmargin=16pt,itemsep=2pt,topsep=3pt,parsep=0pt}
\setlist[enumerate]{leftmargin=20pt,itemsep=5pt,topsep=3pt,parsep=0pt}
\pagestyle{fancy}\fancyhf{}\fancyhead[L]{\small Phurailatpam Hemantakumar}\fancyhead[R]{\small\documenttype}\fancyfoot[R]{\small\thepage}\renewcommand{\headrulewidth}{0.3pt}
\newcommand{\firstheading}[2]{\thispagestyle{plain}{\LARGE\bfseries\color{accent}#1}\par\vspace{3pt}{#2}\par\vspace{4pt}{\color{accent}\hrule}\vspace{7pt}}
\newcommand{\dated}[2]{\par\noindent\begin{tabularx}{\linewidth}{@{}X r@{}}#1&#2\end{tabularx}\par}
\newcommand{\contact}{\href{mailto:hemantaph@link.cuhk.edu.hk}{hemantaph@link.cuhk.edu.hk}\quad|\quad +852 9867 8178\quad|\quad\href{https://hemantaph.com}{hemantaph.com}\\Department of Physics, The Chinese University of Hong Kong}
\setlength{\parskip}{2pt}
\hypersetup{pdftitle={Curriculum vitae | Phurailatpam Hemantakumar}}
\begin{document}
\firstheading{Phurailatpam Hemantakumar}{\contact\\Citizenship: Indian\quad|\quad\href{https://github.com/hemantaph}{github.com/hemantaph}}
\section*{Research interests}
Strong gravitational-wave lensing statistics, compact-binary population simulations, modelling gravitational waves from binaries in eccentric orbits, Bayesian data analysis, and multimessenger lensing.
\section*{Education}
\dated{\textbf{PhD in Physics}, The Chinese University of Hong Kong}{2020--2026}
Thesis: \textit{Statistical and Data Analysis of Eccentric and Lensed Gravitational Waves}.\\
Supervisor: Otto Akseli Hannuksela. Degree award expected December 2026.
\dated{\textbf{M.Sc. in Physics}, Indian Institute of Technology Bombay}{2016--2018}
\dated{\textbf{B.Sc. (Hons.) in Physics}, Hindu College, University of Delhi}{2012--2015}
\section*{Research and scientific software}
\textbf{Strong gravitational-wave lensing statistics.} Lead developer of \href{https://ler.hemantaph.com/}{\texttt{ler}}, a Monte Carlo framework for strongly lensed and unlensed compact-binary populations and detectable rates, with Otto Akseli Hannuksela, Anupreeta More, David Keitel, Chris Van Den Broeck, and other collaborators. Implemented source and lens population models, elliptical power-law lenses with external shear, semi-analytical lens-equation solvers for computing image properties, numerical cross-sections and optical depths, and detector-selection effects. Implemented propagation of population-hyperparameter and finite-sampling uncertainties; analysis and publication are ongoing. Working with Chunglee Kim and Jiyun Park on star-formation-informed merger rates.

\textbf{SNR and detection probability.} Lead developer of \href{https://gwsnr.hemantaph.com/}{\texttt{gwsnr}}, in collaboration with Otto Akseli Hannuksela, using partial-scaling interpolation and parallel waveform inner products for signal-to-noise ratio calculations, with neural-network methods for detectability estimates. Maintain code, documentation, and user support for both packages.

\textbf{Eccentric waveform development and implementation.} With Achamveedu Gopakumar, Maria Haney, Tjonnie Li, and Srishti Tiwari, co-developed \href{https://arxiv.org/abs/2508.12697}{\texttt{TaylorF2Ecck}}, an analytic, frequency-domain post-Newtonian extension of \texttt{TaylorF2} for nonspinning, low-eccentricity inspirals, based on our earlier work, \href{https://arxiv.org/abs/1905.07956}{Tiwari et al. (2019)}, and implemented it in a \texttt{LALSuite} fork. Conducted Bayesian analyses of GW170817 and GW190425 and injection studies of waveform accuracy. In a separate study with Lalit Pathak and Achamveedu Gopakumar, we investigated the degeneracy between \href{https://arxiv.org/abs/2605.21955}{eccentricity and line-of-sight acceleration in GW190814}.

\textbf{Multimessenger lensing.} With Anupreeta More, developed \href{https://arxiv.org/abs/2502.02536}{detectability forecasts for lensed short gamma-ray bursts and binary-neutron-star GW counterparts}. In collaboration with Otto Akseli Hannuksela, Albert Kong, and other collaborators, developing Bayesian posterior-odds tests for multiple GRBs using sky-localisation and time-delay information.

\textbf{Additional experience.} Familiarity with binary population synthesis through a \href{https://docs.google.com/presentation/d/1WBn-nu949xdn9oGZe1GgIlbHsQ_ZBelahJA6aSFyVq0/edit}{journal club presentation} on \href{https://arxiv.org/abs/2405.01623}{Boesky et al. (2024)}, for which I reproduced results using \href{https://compas.science/}{COMPAS}.
\section*{LVK/GRASP commitments and service}
LIGO Scientific Collaboration member. LVK/GRASP commitments: project leadership, development, testing, documentation, and user support for \texttt{ler}, including ISR-weighted work towards GWTC-6 and GWTC-7; continued \texttt{gwsnr} development and support. Planned GRASP activities also include eccentric waveform development and implementation.
\newpage
\textbf{Prospective research.} Explore extending \texttt{ler} statistics to millilensing and microlensing.
\section*{Computational and statistical methods}
\textbf{Programming:} Python, C, C++, Fortran, Mathematica; NumPy, SciPy, Numba, JAX, MLX.\\
\textbf{Methods:} Monte Carlo simulation, importance and rejection sampling, numerical integration, interpolation, Bayesian parameter estimation and model comparison, nested sampling, selection effects, and parallel CPU/GPU computation. \texttt{LALSuite}, \texttt{bilby}, Git, Linux, and \LaTeX.
\section*{Research mentoring, teaching, and outreach}
\dated{\textbf{Jiyun Park}, PhD student, Ewha Womans University}{Oct 2026--present}
Mentoring on star-formation-informed binary-black-hole merger densities and GW lensing rates.
\dated{\textbf{Jihwan Park}, M.Sc. student, Ewha Womans University}{Jan--May 2026}
Mentored a project on data-driven predictions of GW lensing rates.
\dated{\textbf{Yiwei Li}, CUHK SURIP}{May--Jun 2024}
Supervised research on \texttt{ler} statistics and documentation, and galaxy-cluster lensing.

\textbf{Additional mentoring:} Brian Hiu Yeung Cheng (CUHK SURE, Apr--Jun 2024), Laura E. Uronen (early PhD, Aug--Oct 2023), and Leo C. Y. Ng (CUHK SURE, Feb--Jun 2023). Work with \href{https://arxiv.org/abs/2406.14257}{Uronen} and \href{https://arxiv.org/abs/2403.16532}{Ng} led to co-authored publications in 2025.

\textbf{Teaching assistant, CUHK:} General Physics for Engineers (2020); Astronomy (2020--2022); Engineering Physics: Mechanics and Thermodynamics (2021).

\textbf{Workshop instruction:} Served as a resource person and conducted GW data-analysis workshops at Manipur University (Nov 2025 and 20--21 Oct 2026). Taught \href{https://gw.phy.cuhk.edu.hk/activities/surip-summer-school-2026/}{\textit{Python Foundations for Scientific Problem Solving}} at the CUHK SURIP Summer School (Jun 2026); conducted GW sessions and seminars at Dhanamanjuri University (Nov 2025).
\section*{Selected seminars and conference presentations}
\begin{itemize}
\item \textbf{Oct 2026:} \textit{Constraining Gravitational-Wave Events}, seminar, RIST and Manipur University.
\item \textbf{Sep 2026:} \textit{Detectable Strongly Lensed Gravitational Waves: Populations, Rates, and Selection Effects}, seminar, IUCAA, Pune.
\item \textbf{Sep 2026:} \texttt{TaylorF2Ecck}, poster, LVK Collaboration Meeting, Pune.
\item \textbf{Dec 2025:} \texttt{ler}, talk, East-Asia Gravitational-Wave Science Workshop, Nagoya.
\item \textbf{Nov 2022:} Eccentricity and subdominant harmonics in GW parameter estimation, invited talk, 6th Workshop on Gravitational Waves, Seoul National University.
\item \textbf{Jan 2022:} \textit{Nested Sampling: An insight on the Bilby parameter estimation of gravitational waves}, invited talk, GWnext 2022.
\end{itemize}
\section*{Fellowships and awards}
\dated{Postgraduate studentship, Department of Physics, CUHK}{Awarded 2020}
\dated{Karlsruhe House of Young Scientists internship grant}{2017}
\dated{DST INSPIRE Fellowship, Government of India}{2012--2016}
\dated{State Merit Scholarship, Higher Education Manipur}{2012--2015}
\section*{Earlier research experience}
\textbf{AMS-02 internship, Karlsruhe Institute of Technology} (May--Jul 2017). Analysed helium and carbon cosmic-ray events, differential rates, and spatial distributions using HEALPix and C++.
\section*{Publications and products}
The accompanying publication list separates refereed papers, preprints, a white paper, conference proceedings, science writing, and scientific software.
\end{document}
